% !TeX program = xelatex
% 唯一项目正文；图像和原始建模证据在 report-assets/。
\documentclass[UTF8,fontset=fandol,zihao=-4,oneside]{ctexrep}
\usepackage[a4paper,margin=2.2cm]{geometry}
\usepackage{graphicx,longtable,booktabs,array,tabularx,xcolor}
\usepackage{enumitem,fvextra,float,fancyhdr}
\usepackage{xurl}
\usepackage[unicode,hidelinks]{hyperref}
\setlength{\headheight}{16pt}
\pagestyle{fancy}
\fancyhf{}
\fancyhead[L]{校车管理系统：综合项目报告}
\fancyhead[R]{BUS-FLUTTER}
\fancyfoot[C]{\thepage}
\setlist{nosep,leftmargin=2em}
\setlength{\parskip}{0.35em}
\setlength{\emergencystretch}{3em}
\renewcommand{\arraystretch}{1.2}
\DefineVerbatimEnvironment{code}{Verbatim}{breaklines=true,breakanywhere=true,fontsize=\small,frame=single}
\newcommand{\pathcode}[1]{\nolinkurl{#1}}
\newcommand{\toolmark}[1]{\par\noindent\textbf{课程工具标识：}#1\par}
\newcommand{\reportfigure}[2]{%
  \begin{figure}[H]\centering
  \includegraphics[width=\linewidth,height=0.78\textheight,keepaspectratio]{#1}
  \caption{#2}\end{figure}}
\hypersetup{pdftitle={校车管理系统综合项目报告},pdfauthor={BUS-FLUTTER 项目组}}

\begin{document}
\begin{titlepage}
\centering
\vspace*{2cm}
{\Huge\bfseries 校车管理系统\par}
\vspace{0.8cm}
{\LARGE 综合项目报告与技术文档\par}
\vspace{0.6cm}
{\Large Flutter 跨平台三端课程原型\par}
\vspace{2cm}
\begin{tabular}{rl}
课程 & 软件开发环境与工具 \\
仓库 & BUS-FLUTTER \\
工程 & school\_bus\_app \\
主要技术 & Flutter / Dart / Riverpod / GoRouter \\
课程交付计划 & 2026 年 4 月 16 日至 6 月 25 日 \\
文档整合日期 & 2026 年 10 月 10 日 \\
修订日期 & 2026 年 11 月 8 日 \\
文档版本 & 综合修订版 1.1 \\
\end{tabular}
\vfill
本文统一收录项目计划、需求、设计、源码说明、测试、配置管理、\par
工具使用证据和后续技术方案；当前实现与历史记录、规划能力分开表述。\par
\vspace{1cm}
Overleaf 主文件：\texttt{main.tex}\quad 编译器：XeLaTeX\par
\end{titlepage}

\pagenumbering{roman}
\chapter*{摘要与阅读说明}
\addcontentsline{toc}{chapter}{摘要与阅读说明}
校车管理系统面向学生、教职工、司机及校园运营管理人员，采用一个 Flutter 工程提供普通用户、司机、后台管理三类入口。应用以 Mock 样例初始化 Riverpod 状态，演示预约选座、模拟支付、取消扣分、司机任务与模拟定位、车辆司机维护、调度建议确认、位置展示和统计看板。2026 年 11 月修订增加网页 localStorage 保存和同源标签页同步，刷新后可恢复业务数据；原生平台仍使用会话内存。认证会话和密码不保存，当前没有真实后端或跨设备同步。

本文不是把原文件逐份附在末尾，而是将相同内容去重，按软件生命周期重新组织，并保留原有 UML 图、计划截图、后台测试截图及可编辑建模证据。所有正文集中于本文件；辅助文件不构成另一套报告。

\textbf{事实层级：}
\begin{itemize}
\item 当前实现：以 \pathcode{pubspec.yaml}、\pathcode{lib/}、\pathcode{test/}、CI 工作流及本次实际验证为依据。
\item 历史记录：原计划、测试报告、截图中的日期、费用、通过结论仅代表原材料记载，不等于本次重新执行。
\item 扩展方案：REST、WebSocket、Supabase 客户端、真实支付、GPS、FCM、地图服务、OpenAI 等尚未接入的能力单独列入规划章节。
\end{itemize}

\textbf{导入 Overleaf：}选择 New Project → Upload Project，上传 \texttt{overleaf.zip}；将主文件设置为 \texttt{main.tex}，编译器设置为 XeLaTeX，然后 Recompile。中文排版采用 TeX Live 自带 Fandol 字体，不要求上传系统字体；代码使用普通 Verbatim，不依赖 minted、Python 高亮或 shell-escape。

\tableofcontents
\listoffigures
\clearpage
\pagenumbering{arabic}

\chapter{项目概述与实现范围}
\section{背景、目标和角色}
跨校区、宿舍区、教学楼、科技园与附属医院之间存在持续的校园通勤需求。人工预约和电话派车容易造成余座信息不透明、通知滞后、调度缺乏数据依据和运行状态不可见。项目目标是以三类工作台将这些流程系统化，并提供课程所要求的计划、建模、实现、测试及配置管理证据。

\begin{longtable}{p{2cm}p{5.7cm}p{7cm}}
\toprule 角色 & 目标 & 当前功能 \\ \midrule
普通用户 & 查询校车、管理预约与费用 & 登录、浏览车次、选座、预约列表、模拟支付、取消、信用等级、本地通知 \\
司机 & 明确任务并展示运行状态 & 登录、任务卡片、乘客名单与上车点、日周月统计、开始任务、模拟位置记录、到站完成 \\
管理员 & 维护运力并辅助调度 & 车辆和司机维护、绑定校验、人工调度、规则式智能建议、确认生效、模拟位置监控、示例 ETA、历史统计 \\
\bottomrule
\end{longtable}

\section{当前实现与生产能力的区别}
\begin{longtable}{p{3cm}p{5.5cm}p{6.2cm}}
\toprule 能力 & 仓库中的实际实现 & 不应误认为已实现的内容 \\ \midrule
三端入口 & 编译期 \texttt{APP\_MODE} 与页面入口切换 & 三套独立部署且同步数据的业务服务 \\
认证 & 三个演示账号、450ms 模拟延迟与角色匹配 & 注册、真实身份验证、Token 和服务端授权 \\
数据 & Mock 初始化；网页 localStorage 保存业务快照、同源标签页同步；原生为会话内存 & 服务端数据库、跨浏览器/设备同步、事务与并发一致性保证 \\
支付 & 选择微信或支付宝标签，500ms 后更新本地状态 & SDK 下单、扣款、支付回调、真实退款和对账 \\
通知 & \texttt{notificationProvider} 本地通知列表 & FCM 或其他平台推送 \\
定位监控 & 5 秒 Timer 增量模拟经纬度，自绘网格地图 & GPS 权限与采集、后台定位、WebSocket 广播、地图 SDK \\
客户 ETA & 固定的 1.2/2.6/4.8km、8/14/23 分钟示例 & 实际乘客定位、Haversine 计算和路网 ETA \\
智能调度 & 按需求人数降序，匹配空闲且已绑定司机的车辆 & OpenAI 调用、全局最优、多车分拆或真实调度执行 \\
Supabase & 表结构、RLS 与种子 SQL 已保留 & App 中可用的 Supabase 后端切换 \\
\bottomrule
\end{longtable}
管理员确认调度会修改后台自身状态，并增加后台位置记录，但不会创建乘客端车次或司机端任务。因而目前是三种业务流程的课程演示原型，不应称为跨客户端实时闭环系统。

\chapter{项目计划与管理}
\toolmark{Microsoft Project / Project 类工具、CSV 导入表、甘特图、资源与成本计划}
\section{计划口径与时间差异}
原计划说明记载 2026-04-16 至 2026-06-04、36 个工作日、400 工时、29 个任务和预算 551,700 元；CSV 及原最终报告采用 2026-04-16 至 2026-06-25 的课程交付周期。两份资料不能视为相同进度基线。

本报告用 CSV 展示三端原型的课程交付计划，同时保留原 Project 计划中的成本、资源与截图作为历史计划证据。6 月 4 日为原计划结束日期，6 月 25 日为课程提交计划日期；均不据此推断实际完成时间。CSV 中持续时间与开始、结束日期含周末，未经指定 Project 工作日历重算，不据此计算精确关键路径。

\section{课程交付工作分解}
\small
\begin{longtable}{p{0.6cm}p{4cm}p{0.8cm}p{3.8cm}p{1.5cm}p{3cm}}
\toprule ID & 任务 & 工期 & 开始至结束（2026年） & 前置 & 责任资源 \\ \midrule
1 & 项目启动和选题确认 & 2d & 04-16 至 04-17 & — & 项目负责人 \\
2 & 项目范围和计划编制 & 5d & 04-18 至 04-22 & 1 & 项目负责人 \\
3 & 需求调研和角色分析 & 7d & 04-23 至 04-29 & 2 & 需求分析人员 \\
4 & 用例和业务规则建模 & 5d & 04-30 至 05-04 & 3 & 需求分析人员 \\
5 & 需求评审和需求基线 & 2d & 05-05 至 05-06 & 4 & 负责人、需求人员 \\
6 & 概要设计和架构选型 & 4d & 05-07 至 05-10 & 5 & 开发人员 \\
7 & UML 和核心流程设计 & 5d & 05-11 至 05-15 & 6 & 开发人员 \\
8 & 数据模型和接口设计 & 4d & 05-16 至 05-19 & 7 & 开发人员 \\
9 & 普通用户端实现 & 7d & 05-20 至 05-26 & 8 & 开发人员 \\
10 & 司机端实现 & 5d & 05-27 至 05-31 & 8 & 开发人员 \\
11 & 后台车辆和司机维护 & 6d & 06-01 至 06-06 & 8 & 开发人员 \\
12 & 调度、监控、分析实现 & 6d & 06-07 至 06-12 & 11 & 开发人员 \\
13 & 支付信用通知及 UI 收口 & 4d & 06-13 至 06-16 & 9 & 开发人员 \\
14 & 自动化测试设计 & 2d & 06-17 至 06-18 & 9, 10, 12, 13 & 测试人员 \\
15 & 测试执行和缺陷修复 & 3d & 06-19 至 06-21 & 14 & 测试、开发人员 \\
16 & 配置管理和交付基线 & 1d & 06-22 & 15 & 配置管理员 \\
17 & 报告整理和证据汇总 & 3d & 06-22 至 06-24 & 15 & 项目负责人 \\
18 & 最终检查和提交 & 1d & 06-25 & 16,17 & 项目负责人 \\
\bottomrule
\end{longtable}
\normalsize
可编辑证据位于 \pathcode{report-assets/evidence/project_plan.csv} 和 \pathcode{report-assets/evidence/project_plan.mpp}。CSV 是上述课程交付计划，MPP 是原 Project 计划；导入时需分别核对日历、任务依赖与资源，不能声称两者完全一致。

\section{原 Project 计划与关键路径}
\begin{longtable}{p{3cm}p{2cm}p{4.2cm}p{5.5cm}}
\toprule 阶段 & 原计划工期 & 原计划区间 & 工作内容 \\ \midrule
项目启动 & 8 工作日 & 04-16 至 04-27 & 需求分析、可行性研究 \\
系统设计 & 6 工作日 & 04-28 至 05-05 & 概要设计、详细设计 \\
编码实现 & 11 工作日 & 05-06 至 05-20 & 数据库、前端、后端逻辑、接口 \\
测试 & 8 工作日 & 05-21 至 06-01 & 单元、集成、系统测试 \\
部署上线 & 3 工作日 & 06-02 至 06-04 & 环境部署、项目验收 \\
沟通管理 & 30.25 工作日 & 04-20 至 06-01 & 周会、评审、阶段汇报 \\
\bottomrule
\end{longtable}
原材料标注的阶段关键路径为“项目启动 → 系统设计 → 编码实现 → 测试 → 部署上线”，并列出可行性研究、详细设计、后端逻辑开发、接口开发、系统测试、环境部署为关键任务。这是原计划的描述，不是本次对 MPP 重新计算的结果。
原里程碑为启动完成 4 月 27 日、设计完成 5 月 5 日、编码完成 5 月 20 日、测试完成 6 月 1 日、验收 6 月 3 日。
\reportfigure{report-assets/plan/plan_image1.png}{原 Project 工作分解与任务表：历史计划截图}
\reportfigure{report-assets/plan/plan_image2.png}{原 Project 甘特图、依赖和资源分配：历史计划截图}

\section{资源和成本}
\begin{longtable}{p{3.7cm}p{3cm}p{3cm}p{5cm}}
\toprule 人力资源 & 标准费率（元/工时） & 加班费率（元/工时） & 原计划最大单位 \\ \midrule
项目经理 & 500 & 750 & 100\% \\
系统分析师 & 400 & 600 & 100\% \\
前端工程师 & 300 & 450 & 100\% \\
后端工程师 & 300 & 450 & 100\% \\
测试工程师 & 250 & 375 & 100\% \\
UI 设计师 & 350 & 525 & 100\% \\
\bottomrule
\end{longtable}
材料单价为开发服务器 20,000 元/台、测试服务器 15,000 元/台、办公电脑 8,000 元/台、网线耗材 5 元/米。它们是原计划模拟预算，并非真实采购或实际项目支出。

\begin{longtable}{p{3cm}p{7cm}p{4.7cm}}
\toprule 任务 & 资源分配 & 数量或比例 \\ \midrule
需求分析 & 系统分析师、办公电脑 & 100\%、2 台 \\
可行性研究 & 分析师、经理、电脑 & 各 50\%、2 台 \\
概要设计 & 分析师、电脑 & 100\%、3 台 \\
详细设计 & UI 设计师、分析师、电脑 & 各 100\%、3 台 \\
数据库 / 前端 / 后端 & 对应工程师、开发服务器、电脑 & 工程师 100\%、1 台、2 台 \\
接口开发 & 前后端工程师、服务器、电脑 & 各 50\%、1 台、2 台 \\
单元 / 集成测试 & 测试工程师、测试服务器、电脑 & 100\%、1 台、2 台 \\
系统测试 & 测试工程师、经理、测试服务器、电脑 & 各 50\%、1 台、2 台 \\
环境部署 & 后端工程师、开发服务器、测试服务器、电脑、网线 & 100\%、2 台、1 台、5 台、100 米 \\
项目验收 & 项目经理 & 100\% \\
\bottomrule
\end{longtable}
原材料记录一个过度分配的工时资源，但未明确冲突人员与解决结果；后续重排计划需检查资源平衡。

\begin{longtable}{p{3.7cm}r p{6.5cm}}
\toprule 阶段 & 预算（元） & 子任务预算（元） \\ \midrule
项目启动 & 52,800 & 需求分析 27,600；可行性研究 23,200 \\
系统设计 & 80,400 & 概要设计 31,400；详细设计 48,000 \\
编码实现 & 196,800 & 数据库 48,000；前端 50,400；后端 55,200；接口 43,200 \\
测试 & 111,000 & 单元 39,000；集成 35,000；系统 37,000 \\
部署上线 & 110,700 & 环境部署 102,700；验收 8,000 \\
总计 & 551,700 & 原计划总预算，不含真实结算信息 \\
\bottomrule
\end{longtable}
需求分析另记固定成本 2,000 元，概要设计另记 1,000 元；其他列出的任务固定成本为 0。阶段汇总与子任务金额存在差额，原表未完整说明计算口径，不能简单相加后称为已审计预算。编码阶段占原总预算约 35.7\%。现金流峰值记录为 6 月 1 日当周 129,200 元；完整历史现金流图如下。
\reportfigure{report-assets/plan/plan_image3.png}{原计划按周现金流：保留原预算分布，不作为真实财务记录}

\section{沟通、质量和风险管理}
每周项目组会议形成纪要；阶段节点开展需求评审并形成评审报告；里程碑向项目负责人汇报并形成阶段报告。
原计划的交叉评审方式可保留，但其中“Java 开发手册”不适用于该 Flutter 工程；当前规范以 \pathcode{analysis_options.yaml} 引入的 \texttt{flutter\_lints} 为准。原覆盖率不低于 80\% 是计划目标，没有现成覆盖率证据证明达标。

\begin{longtable}{p{2cm}p{4.2cm}p{2cm}p{6.5cm}}
\toprule 风险 & 描述 & 原概率 / 影响 & 应对措施 \\ \midrule
技术 & 新技术学习成本 & 中 / 高 & 提前培训、先做原型、保留缓冲时间 \\
人员 & 核心人员离开 & 低 / 高 & 文档化、版本管理和知识共享 \\
进度 & 需求频繁变更 & 高 / 中 & 评估范围、记录需求变更、同步测试 \\
验证 & 模拟数据被误称为生产能力 & 已知限制 & 区分原型、历史证据与未来服务；不以截图替代实测 \\
\bottomrule
\end{longtable}

\chapter{需求规格与验收}
\toolmark{StarUML、需求规格说明、需求追踪矩阵}
\section{业务需求和用例}
业务目标包括提升预约效率（BRQ-01）、提高任务透明度（BRQ-02）、支持调度决策（BRQ-03）、掌握运行状态（BRQ-04）、提供运营复盘（BRQ-05）、满足课程交付（BRQ-06）。其中运行状态和运营分析当前采用演示数据。

\begin{longtable}{p{1.4cm}p{4cm}p{2.2cm}p{6.8cm}}
\toprule 编号 & 用例 & 参与者 & 触发与输出 \\ \midrule
UC-01 & 登录系统 & 全部角色 & 账号密码与入口匹配后进入工作台 \\
UC-02 & 查看车次 & 普通用户 & 展示路线、时间、票价、余座 \\
UC-03 & 预约车次 & 普通用户 & 选择可用座位、创建预约 \\
UC-04 & 支付车费 & 普通用户 & 待支付订单选择方式、模拟完成 \\
UC-05 & 取消预约 & 普通用户 & 取消待发车预约、记录信用扣分 \\
UC-06 & 查看信用等级 & 普通用户 & 展示信用分、等级和取消记录 \\
UC-07 & 查看通知 & 普通用户 & 查看预约、支付、信用通知 \\
UC-08 & 查看派车任务 & 司机 & 查看今日任务卡片 \\
UC-09 & 查看乘客名单 & 司机 & 展示姓名、电话尾号、上车点、座位 \\
UC-10 & 日周月统计 & 司机 & 切换任务、人次、里程、准点率 \\
UC-11 & 共享位置 & 司机 & 开始任务后追加模拟位置记录 \\
UC-12 & 车辆管理 & 管理员 & 新增、编辑、删除车辆 \\
UC-13 & 司机管理 & 管理员 & 维护司机档案与车辆绑定 \\
UC-14 & 人工调度 & 管理员 & 按需求选择车辆、司机并生成方案 \\
UC-15 & 智能辅助调度 & 管理员 & 本地规则生成推荐，人工确认 \\
UC-16 & 位置监控 & 管理员 & 选择模拟标记查看车辆、司机、速度 \\
UC-17 & 查看客户 ETA & 管理员 & 展示固定客户距离和时间示例 \\
UC-18 & 历史数据分析 & 管理员 & 切换周期查看人次、收入、利用率、热度 \\
\bottomrule
\end{longtable}
原 UML 使用 include/extend 表达预约与查询、支付与预约、通知与取消、智能建议与人工调度的关系。它们描述设计意图，不代表后台与司机端已经连接实时服务。

\section{功能需求}
\begin{longtable}{p{1.2cm}p{3.5cm}p{5cm}p{4.7cm}}
\toprule 编号 & 功能 & 输入及处理 & 输出 / 当前范围 \\ \midrule
FR-01 & 登录与角色入口 & 手机、密码、入口模式，校验账号和角色 & 工作台或错误提示；无注册 \\
FR-02 & 预约选座 & 用户、车次、座位，检查时刻、余座、重复 & 待发车且待支付的预约 \\
FR-03 & 预约查看筛选 & 用户 ID、预约状态 & 对应订单列表 \\
FR-04 & 取消与信用扣分 & 预约 ID，检查状态与发车时间 & 已取消、扣分记录、通知 \\
FR-05 & 车费支付 & 预约 ID、模拟支付方式 & 支付记录与本地支付通知 \\
FR-06 & 信用与通知 & 信用分、取消记录、通知状态 & 信用面板与本地通知中心 \\
FR-07 & 任务及乘客名单 & 仓库中的司机任务 & 路线、车次、上车点和座位 \\
FR-08 & 司机统计 & 日、周、月 & 完成车次、人次、里程、准点率 \\
FR-09 & 位置共享 & 开始 / 完成任务 & 每 5 秒模拟位置记录，结束后停止 \\
FR-10 & 车辆管理 & 车牌、车型、座位数、状态 & 更新列表及校验错误 \\
FR-11 & 司机管理 & 姓名、手机、驾照、车辆绑定 & 档案列表与一对一绑定校验 \\
FR-12 & 人工与智能调度 & 待调度需求、空闲车辆、司机 & 待确认建议及确认后的后台状态 \\
FR-13 & 位置与 ETA & 后台模拟位置、固定客户样例 & 自绘地图标记和示例 ETA \\
FR-14 & 运营分析 & 日、周、月统计周期 & 演示趋势、人次、收入、利用率、热度 \\
\bottomrule
\end{longtable}
FR-06、FR-08、FR-14 为原文中等优先级，其余为高优先级。真实地图、乘客位置上传、CSV 数据导出、自定义统计时间段属于扩展需求，当前并未全部实现。

\section{业务规则与状态语义}
\begin{longtable}{p{1.4cm}p{8.5cm}p{4.7cm}}
\toprule 规则 & 当前约束 & 关联模块 \\ \midrule
RUL-01 & 账号角色必须与当前入口一致 & MockAuthRepository、路由 \\
RUL-02 & 只有 scheduled 且未过发车时刻的车次可预约 & TripModel.isBookable \\
RUL-03 & 同用户同车次不能存在重复 active 预约；座位必须可选 & \pathcode{PassengerBookingController} \\
RUL-04 & 只有 pending 且未发车的预约可取消 & cancelBooking \\
RUL-05 & 距发车整数分钟数不大于 120 扣 5 分，否则扣 2 分；信用分限制 0--100 & 取消规则、AuthController \\
RUL-06 & UI 与 Controller 仅允许 pending 且 unpaid 进入新付款；processing 拒绝重复提交 & BookingModel.canPay、payBooking \\
RUL-07 & 运行中车辆不可删除 & \pathcode{AdminController.deleteVehicle} \\
RUL-08 & 有后台位置记录的司机不可删除 & \pathcode{AdminController.deleteDriver} \\
RUL-09 & 车辆与司机保持一对一绑定；司机手机不能重复 & 驾驶员维护校验 \\
RUL-10 & 生成和确认均校验未分配且未到发车时间的需求、空闲车辆、足够座位、司机可用与双向绑定 & 调度统一校验 \\
RUL-11 & 智能调度仅选空闲且已绑定司机车辆，同批不重复用车 & \pathcode{generateAiDispatchPlans} \\
RUL-12 & 一名司机的本地任务状态只能有一个 activeTaskId & DriverTaskController \\
\bottomrule
\end{longtable}
Controller 已拦截同一订单正在处理时的重复支付，并在延迟回调时核对当前付款尝试和订单状态。这是本地演示逻辑，正式后端仍需独立实现授权、幂等和事务。已支付预约仍是 \texttt{BookingStatus.pending}，支付结果由独立的 \texttt{BookingPaymentStatus.paid} 表示；不能使用原方案中的 \texttt{BookingStatus.paid} 代替。

\section{非功能需求与验证方式}
\begin{longtable}{p{1.2cm}p{2.5cm}p{6.5cm}p{4.5cm}}
\toprule 编号 & 类别 & 要求 & 证据或限制 \\ \midrule
NFR-01 & 易用性 & 卡片式工作台、明确的核心按钮 & Widget 测试、历史界面截图 \\
NFR-02 & 可维护性 & core/features/shared 分层，业务放 Notifier & 源码结构核对 \\
NFR-03 & 可测试性 & 核心主流程和负向业务规则可自动验证 & Flutter Test \\
NFR-04 & 可扩展性 & 数据初始化与页面分离 & Mock 仓库；真实仓库尚未接入 \\
NFR-05 & 响应式 & LayoutBuilder 适配布局 & 当前页面有实现，仍需多设备验证 \\
NFR-06 & 时序 & 5 秒模拟定位、500ms 模拟支付 & 代码常量，不是性能保证 \\
NFR-07 & 安全 & 正式版需加密、权限、审计、隐私保护 & 原型不提供生产安全保证 \\
\bottomrule
\end{longtable}

\section{需求追踪与验收标准}
\begin{longtable}{p{2cm}p{5.6cm}p{7.1cm}}
\toprule 需求 & 主要实现位置 & 测试对应 \\ \midrule
FR-01 & auth、core/providers/auth、core/router & 三种登录 Widget；role entrance mismatch \\
FR-02--03 & passenger/booking、passenger\_home\_page & 选座生成订单；duplicate booking \\
FR-04 & cancelBooking、adjustCreditScore & 信用 96→91；completed/cancelled 拒绝取消 \\
FR-05--06 & payBooking、notificationProvider & 支付成功与本地通知；不代表真实扣款 \\
FR-07--09 & driver/tasks、driver\_home\_page & 司机专项 57 项与主流程 Widget \\
FR-10--11 & adminProvider、车辆与司机模型 & 编辑、新增；运行车辆删除、绑定冲突 \\
FR-12 & dispatch 模型、AdminController & 生成确认；座位不足、无可用车辆 \\
FR-13 & admin\_home\_page、location 模型 & 标记点击、模拟刷新、固定 ETA 展示 \\
FR-14 & analytics 模型与后台统计面板 & 周统计切换与路线热度展示 \\
\bottomrule
\end{longtable}
验收流程分别为乘客“登录→浏览→预约→支付或取消”、司机“登录→查看任务→开始→模拟定位→完成”、管理员“维护运力→生成建议→确认→查看监控”。UML 验收按实际保留的 12 张图计算，而非原需求验收表误写的 10 张。文档验收改为本综合报告完整覆盖七类课程内容，不再要求七份重复文件。测试通过仅说明被覆盖场景，不能等同全部需求已经被充分验证。

\chapter{总体架构与详细设计}
\toolmark{Flutter / Dart、Riverpod、GoRouter、StarUML}
\section{真实依赖与运行入口}
\begin{longtable}{p{4.5cm}p{3.5cm}p{6.7cm}}
\toprule 项目 & 声明或锁定版本 & 用途 \\ \midrule
Dart & SDK 约束 \texttt{\^{}3.11.5} & 语言与工具链要求 \\
Flutter & CI 固定 3.47.9 stable & 六平台工程；不代表六平台均经测试 \\
flutter\_riverpod & 声明及锁定 3.3.1 & Provider、Notifier 与状态订阅 \\
go\_router & 声明及锁定 17.2.2 & 登录跳转与角色路由 \\
cupertino\_icons & 声明 1.0.8，原锁定 1.0.9 & 图标 \\
flutter\_lints & 6.0.0 & 分析规则 \\
flutter\_test & Flutter SDK 内置 & Widget、业务状态与模型测试 \\
web & 1.1.1 & 网页 localStorage 与 storage 事件互操作 \\
\bottomrule
\end{longtable}
当前未引入 Dio、Retrofit、Hive、secure storage、socket\_io\_client、geolocator、flutter\_map、fl\_chart、pay、firebase\_messaging、supabase\_flutter、freezed 或 build\_runner；这些仅是初期选型候选。

\begin{code}
// lib/main.dart: actual initialization pattern
const rawMode = String.fromEnvironment(
  'APP_MODE', defaultValue: 'passenger');
final initialMode = AppMode.fromValue(rawMode);
runApp(ProviderScope(
  overrides: [
    initialAppModeProvider.overrideWithValue(initialMode),
    demoPersistenceProvider.overrideWith((ref) {
      final store = createDemoPersistence();
      ref.onDispose(store.dispose);
      return store;
    }),
  ],
  child: const AppRoot(),
));
\end{code}
\pathcode{lib/app.dart} 观察入口和路由 Provider，构建 \texttt{MaterialApp.router}。\pathcode{lib/core/router/app_router.dart} 保持同一 GoRouter 实例，监听身份、角色与入口变化后刷新重定向，并在销毁时释放路由。未登录或角色不匹配返回登录页；已登录且路径不属于本角色时跳转到对应工作台。信用分变化不再重建导航，快捷账号跨入口切换不会丢失已填手机号。\texttt{APP\_MODE} 是 dart-define 编译参数，不是运行时读取的 shell 环境变量。

\section{模块边界和数据流}
\begin{longtable}{p{3.7cm}p{5.5cm}p{5.5cm}}
\toprule 边界 & 关键实现 & 职责 \\ \midrule
入口与路由 & main.dart、app.dart、appRouterProvider & 初始化、主题、登录态和角色导航 \\
认证 & AuthRepository 接口、MockAuthRepository、AuthController & 演示账号校验、登录退出、信用更新 \\
乘客 & MockPassengerRepository、PassengerBookingController & 车次、预约、支付、取消规则 \\
司机 & MockDriverRepository、DriverTaskController & 任务、统计、位置 Timer 生命周期 \\
后台 & MockAdminRepository、AdminController & 车辆司机、需求、建议、监控、分析 \\
共享 & core/models、notificationProvider、shared/widgets & 值对象、通知、入口切换、公共卡片 \\
\bottomrule
\end{longtable}
数据流是“页面触发操作 → Notifier 校验与生成新模型 → 更新 state → 页面重建，并保存业务快照”。Mock 仓库提供初始样例，后续状态通过 DemoPersistence 保存；网页从其他同源标签页收到 storage 事件后重新读取对应业务域。认证定义了可替换接口，其他三个模块仍直接依赖对应 Mock 类型；当前没有真实后端或统一 \texttt{BackendBootstrap} 服务。

\section{演示数据保存、同步与重置}
\pathcode{lib/core/persistence/} 定义统一的读取、写入、变化通知及清除接口。网页入口选择 localStorage；原生平台、默认测试容器和不允许浏览器存储的环境使用内存回退。\texttt{passenger} 保存车次、预约和模拟付款，\texttt{driver} 保存任务与有限位置历史，\texttt{admin} 保存后台业务状态，\texttt{notifications} 保存本地通知，\texttt{scores} 只保存用户信用分。登录会话、密码和真实凭据不进入这些快照；刷新后需要重新登录。

同一浏览器中访问相同源的标签页会监听对应业务数据的变化，刷新页面也会读取已存样例。此机制不连接服务端，不跨设备或浏览器同步，不保证并发写入的数据库事务语义，也不会将后台方案自动转换为乘客车次或司机任务。浏览器数据可以被用户清除，因此不能用于真实账务或可信授权。

登录页显示保存范围，并提供带确认的“重置演示数据”。确认后清除当前浏览器或会话的修改，各业务域恢复初始样例并退出登录；取消确认不改数据。演示车次超过发车时间、需要恢复初始座位或重新演示时可使用此入口。信用分由仓库和持久化快照维护，正常退出重登不会抹去取消扣分。

\section{核心数据字典}
以下模型位于 \pathcode{lib/core/models/}；字段为核心摘要，以具体 Dart 定义为最终依据。
\small
\begin{longtable}{p{4.2cm}p{10.5cm}}
\toprule 模型 & 职责及主要字段 \\ \midrule
UserModel & 用户身份与信用：id、name、phone、creditScore、role \\
TripModel & 路线与座位：id、routeName、origin、destination、vehicleId、driverId、departureTime、totalSeats、bookedSeats、fare、status；isBookable 为计算属性 \\
BookingModel & 预约：id、userId、tripId、seatNo、status、fare、createdAt、cancelledAt、creditPenalty、paymentStatus、paymentMethod \\
PaymentModel & 支付记录：id、bookingId、amount、method、status、createdAt、completedAt \\
DriverTaskModel & 任务：id、tripNo、routeName、起终点、vehicleId、plateNo、departureTime、distanceKm、status、passengers \\
PassengerPickup & 上车信息：name、phoneSuffix、pickupPoint、seatNo \\
DriverStats & 统计：period、totalTrips、totalPassengers、totalDistance、onTimeRate \\
LocationUpdate & 司机模拟位置：vehicleId、lat、lng、speed、timestamp \\
VehicleModel & 车辆：id、plateNo、model、seatCount、status \\
AdminDriverModel & 档案：id、name、phone、licenseNo、boundVehicleId \\
DispatchDemandModel & 需求：id、routeName、origin、destination、passengerCount、departureTime、status \\
DispatchPlanModel & 建议：id、demandId、routeName、vehicleId、driverId、passengerCount、departureTime、loadRate、isAiGenerated、isConfirmed \\
VehicleLocationModel & 后台位置：vehicleId、tripNo、driverId、lat、lng、speed、updatedAt \\
AdminAnalyticsSnapshot & 运营统计：period、totalPassengers、totalRevenue、vehicleUtilization、趋势和路线排名 \\
AppNotificationModel & 本地通知：标识、标题、正文、类型、时间、已读状态 \\
\bottomrule
\end{longtable}
\normalsize
用户和车次分别与多个预约关联；预约与支付记录通过 bookingId 关联；司机任务包含乘客上车信息及位置记录；司机可绑定一辆车；需求与建议、建议与车辆和司机通过 ID 关联。模型关系不意味着这些模块的演示数据会自动同步。

\section{枚举与状态转换}
\begin{longtable}{p{4.5cm}p{10.2cm}}
\toprule 当前类型 & 字面值 \\ \midrule
UserRole / AppMode & passenger、driver、admin \\
TripStatus & scheduled、running、completed、cancelled \\
BookingStatus & pending、completed、cancelled \\
BookingPaymentStatus & unpaid、processing、paid、refunded \\
PaymentMethod & wechat、alipay \\
PaymentStatus & processing、paid、failed、refunded \\
DriverTaskStatus & pending、running、completed \\
VehicleStatus & idle、running、maintenance \\
DispatchDemandStatus & pending、assigned \\
AnalyticsPeriod & day、week、month \\
\bottomrule
\end{longtable}
原 UML 中部分枚举来自早期方案，例如预约 paid、任务 cancelled、分析 daily/weekly/monthly。保留图像用于展示原建模过程，但解释当前状态必须按上表和源码。

预约创建后为 pending + unpaid；模拟支付经 processing 到 paid，但预约仍为 pending。取消已支付预约时，预约支付状态与对应模拟流水同步变为 refunded；支付中取消则终止该尝试，将流水置 failed、预约支付状态恢复 unpaid，预约本身为 cancelled。这里没有真实退款网关或支付对账。司机任务 pending → running → completed，通过 activeTaskId 限制同时进行的任务；司机任务不会自动推动乘客车次状态。

\section{乘客流程}
\begin{enumerate}
\item 浏览车次：判断发车时刻和状态。剩余座位为总座位减仓库预占 bookedSeats，再减新 active 预约数；可选座位从 bookedSeats+1 起排除已预约座位。
\item 创建预约：依次检查可预约、未满座、无同用户重复预约、座位仍可用；新记录插入列表并产生本地预约通知。
\item 支付：拒绝 processing 状态的重复请求；生成 processing 流水并标记订单支付中，延迟 500ms 后检查容器仍有效、订单未取消且付款尝试仍匹配，再标记完成和推送本地通知。
\item 取消：检查预约存在、pending、未发车；计算取消扣分，更新状态、当前登录用户信用及本地通知。信用分限制在 0--100。
\end{enumerate}
本轮同时修复支付中取消、模拟退款账本不一致及重复提交；恢复快照时还会处理超时未结束的付款尝试。真实系统仍需订单归属校验、服务端幂等键、并发占座事务、网关回调和真实退款。目前 Controller 方法不是生产 API，本地存储同步也不能替代这些约束。

\section{司机任务与定位流程}
\texttt{startTask} 检查 activeTaskId 和目标任务，标记 running 后立即追加第一条位置，并启动 5 秒周期 Timer。位置从纬度 31.2304、经度 121.4737 开始按记录序号递增，速度也由固定规则生成。
\texttt{simulateNextLocationReport} 提供手动模拟按钮；新位置前插，\texttt{latestLocation} 返回第一项，历史最多保留 200 条。\texttt{completeActiveTask} 标记 completed、取消 Timer 并清除 activeTaskId。退出登录、切换至其他角色或 Provider 销毁同样停止定位 Timer；重新以司机身份登录后，可恢复未完成任务的模拟上报。任务状态本身保留，不以退出登录代替到站完成。
司机统计是仓库提供的日周月样例，不是对所有实时任务持续聚合后的数据。

\section{后台管理、智能调度和分析}
车辆新增检查车牌、车型、座位数及重复车牌；编辑检查目标存在及车牌冲突；删除运行车辆被拦截，删除其他车辆时清理司机绑定和后台位置。司机新增检查完整性、重复手机及车辆一对一绑定，后台位置列表中出现的司机不允许删除。

人工调度校验待处理需求、空闲车辆、司机及容量。智能建议按人数降序处理需求，对每项选择当前候选列表中第一辆未使用且容量足够的车，候选必须为空闲且已有绑定司机；无法覆盖的需求被跳过。这个算法并不拆分单个大需求到多辆车，也不搜索全局最优解。生成建议仅创建待确认模型。

创建与确认方案共用校验：需求仍待调度且未到发车时间，车辆空闲且容量足够，司机没有在途任务，司机与车辆的双向绑定不冲突。确认还检查方案所依据的人数、路线与时间未变，并按当前容量计算满载率；车辆缩座或司机绑定改变后，旧建议不能绕过检查。校验通过后才标记方案已确认、需求 assigned、车辆 running，必要时绑定司机并创建后台模拟位置。位置 Timer 每 5 秒仅更新运行车辆的坐标。地图使用 \texttt{CustomPainter} 绘制网格，统计图表也使用 Flutter 自身组件，并未依赖真实地图或 fl\_chart。

\chapter{UML 建模与图文证据}
\toolmark{StarUML 原生模型、PNG 导出、设计过程截图}
\section{模型文件和适用范围}
可编辑 StarUML 原文件在 \pathcode{report-assets/evidence/BUS-FLUTTER-UML.mdj}。实际保留 12 张 UML 图：用例 2、类图 2、状态图 3、时序图 2、部署图 1、组件图 2。历史资料称使用 StarUML 6.3.0；截图可见试用版界面，本文不再沿用原文“开源软件”的错误表述。
图中的外部支付、推送、数据库和 WebSocket 为目标部署示意；\texttt{PaymentService}、\texttt{BackendBootstrap} 等名称也是设计抽象，不应据此寻找当前仓库中不存在的运行服务。

\section{用例与领域模型}
\reportfigure{report-assets/uml/usecase_overview.png}{系统总览用例图：普通用户、司机与管理员}
\reportfigure{report-assets/uml/usecase_passenger.png}{乘客端用例详图：查询、选座、支付、取消与通知}
\reportfigure{report-assets/uml/class_core.png}{历史核心领域类图：当前字段和状态以第四章修订说明为准}
\reportfigure{report-assets/uml/class_enums.png}{历史枚举与业务类型图：包含早期方案的状态命名}

\section{生命周期、交互与部署}
\reportfigure{report-assets/uml/state_booking.png}{预约状态机设计意图：当前预约状态与支付状态分开存储}
\reportfigure{report-assets/uml/state_trip.png}{车次状态机设计：scheduled、running、completed、cancelled}
\reportfigure{report-assets/uml/state_drivertask.png}{司机任务状态机设计：当前实现无 cancelled 枚举}
\reportfigure{report-assets/uml/sequence_driver.png}{司机位置上报时序设计：页面、Controller 与 Timer}
\reportfigure{report-assets/uml/sequence_passenger.png}{乘客预约支付时序设计：当前支付逻辑直接位于 Controller}
\reportfigure{report-assets/uml/deployment.png}{目标部署图：外部认证、实时通信、支付和推送尚未接入}
\reportfigure{report-assets/uml/component_frontend.png}{前端组件设计：入口、路由、三端工作台与状态层}
\reportfigure{report-assets/uml/component_backend.png}{目标数据层组件图：当前为 Mock 样例、业务状态与网页本地存储，尚无服务端}

\section{StarUML 设计过程}
\begin{enumerate}
\item 用例图：创建 Use Case Diagram，添加三种 Actor 及用例，使用 Association 关联角色与操作，Dependency 标注 include/extend；乘客图只保留乘客相关操作。
\item 类图：按核心数据字典添加 Class、属性与类型，画关联与多重性；另图使用 Enumeration 记录业务状态。
\item 状态图：创建 Statechart Diagram，以 InitialState、State、FinalState 和 Transition 表达对象生命周期；预约、车次、司机任务分别建模。
\item 时序图：建立 Actor、页面、Controller、Timer 等 Lifeline，连接 Message、自调用并核对 sequenceNumber；支付服务生命线属于目标抽象。
\item 部署图：使用 Node 展示客户端、应用层及第三方服务，用通信路径标注 HTTPS、REST、WSS；区分现状与规划。
\item 组件图：按表现、状态、仓库、模型和入口分层，以 Dependency 表示模块依赖；当真实后端接入后再校正图中的预备组件。
\end{enumerate}
图片通过 File → Export Diagram → PNG 导出；菜单支持时可以批量导出。更新时保留既有文件名，并在本正文中修改对应解释，避免恢复独立建模说明文件。
\reportfigure{report-assets/uml/process_staruml_main.png}{历史 StarUML 软件主界面截图}
\reportfigure{report-assets/uml/process_usecase.png}{历史用例图设计过程截图}
\reportfigure{report-assets/uml/process_class.png}{历史领域类图设计过程截图}
\reportfigure{report-assets/uml/process_class_enums.png}{历史枚举类图设计过程截图}
\reportfigure{report-assets/uml/process_state.png}{历史状态图设计过程截图}
\reportfigure{report-assets/uml/process_sequence.png}{历史时序图设计过程截图}
\reportfigure{report-assets/uml/process_deployment.png}{历史部署图设计过程截图}
\reportfigure{report-assets/uml/process_component.png}{历史组件图设计过程截图}

\chapter{平台工程与 Windows 宿主设计}
\toolmark{Flutter 平台工程、Win32 C++、CMake}
\section{平台边界}
工程包含 Android、iOS、Web、Linux、macOS、Windows 的 Flutter 平台目录。业务入口仍是 Dart 的 \pathcode{lib/main.dart}，平台 runner 只是宿主，不能把原 Windows 设计报告中的“考勤、设备告警、日志管理”等目标当作已实现业务。
Android 的 Java / Kotlin 编译目标为 17，Gradle 配置采用 AGP 8.11.1、Kotlin 2.2.20；release 仍使用 debug 签名和示例 applicationId，不能直接称为正式上架配置。Linux runner 使用 GTK 3，其 CMake 安装后的 bundle 才包含完整资源，不能运行 intermediates\_do\_not\_run 中的未打包程序。
iOS 启动图可替换 \pathcode{ios/Runner/Assets.xcassets/LaunchImage.imageset/} 中的图片，或在 macOS 打开 \pathcode{ios/Runner.xcworkspace} 后在 Xcode 的 Runner/Assets.xcassets 修改。这一说明整合自平台模板 README。

\section{Windows 入口、窗口和生命周期}
\begin{longtable}{p{4.2cm}p{10.5cm}}
\toprule 文件 / 类 & 责任与关键方法 \\ \midrule
main.cpp / wWinMain & 附加控制台、初始化 COM、配置 data 目录 DartProject、解析入口参数、创建 1280×720 的 school\_bus\_app 窗口、运行消息循环、反初始化 COM \\
Win32Window & 封装 HWND、子窗口、Point、Size；Create、Show、Destroy、SetChildContent、GetClientArea、MessageHandler、UpdateTheme \\
WindowClassRegistrar & 注册和注销原生窗口类，避免重复注册 \\
FlutterWindow & 继承 Win32Window，保存 DartProject，使用 unique\_ptr 管理 FlutterViewController \\
\pathcode{FlutterWindow.OnCreate} & 创建并检查引擎和视图、RegisterPlugins、挂载子视图、首帧显示和 ForceRedraw \\
\pathcode{FlutterWindow.OnDestroy} & 释放 Flutter 控制器，再进行父类销毁 \\
\pathcode{FlutterWindow.MessageHandler} & 优先转发引擎与插件，再处理字体变化并交给父类 \\
utils.h / utils.cpp & CreateAndAttachConsole、GetCommandLineArguments、Utf8FromUtf16 \\
resource.h / Runner.rc & 图标和版本资源，图标 ID 为 101 \\
\pathcode{generated_plugin_registrant} & Flutter 生成的插件注册入口，不手动编辑 \\
\bottomrule
\end{longtable}
启动顺序为 wWinMain → COM / 参数 → DartProject → FlutterWindow / Win32Window → FlutterViewController → 插件注册 → 子窗口挂载 → 首帧 → Win32 消息循环。退出时销毁视图、原生窗口与 COM。消息处理优先级为 Flutter 引擎和插件、自定义 FlutterWindow、Win32Window、默认窗口过程。

\section{构建、容错与扩展}
\pathcode{windows/runner/CMakeLists.txt} 要求 CMake 3.14，创建 WIN32 可执行目标；链接 flutter、flutter\_wrapper\_app、dwmapi.lib；注入版本宏并定义 NOMINMAX；依赖 \texttt{flutter\_assemble}。新增原生源码需加入 runner 源文件列表；应用可执行名称由顶层 \pathcode{windows/CMakeLists.txt} 的 BINARY\_NAME 控制，不应只修改 runner 子目标名。

在 Windows、具备匹配 Flutter 要求的 Visual Studio C++ 桌面开发工具链时使用：
\begin{code}
flutter run -d windows --dart-define=APP_MODE=admin
flutter build windows --release --dart-define=APP_MODE=admin
\end{code}
发行目录需整体分发，包含 EXE、Flutter DLL、插件和 data，不能只拷贝一个 EXE。当前 BINARY\_NAME 为 school\_bus\_app，原设计文档“Runner.exe”只是泛称；应以实际产物为准。本次 Linux 环境未执行 Windows、Android 或 iOS 构建。

窗口创建失败和引擎无效时返回失败，智能指针管理视图控制器生命周期；字符编码转换检查空指针；DPI 和主题处理采用 Windows 平台实现。具体兼容性范围需要在目标 Windows 版本实测，而不是引用历史设计的版本声明作为保证。Platform Channel、串口/GPS/门禁、多窗口和本地日志属于后续扩展，当前没有相应业务通道。

\chapter{源码使用、运行与演示}
\toolmark{Flutter / Dart、编辑器、Git}
\section{开发命令与账号}
\begin{code}
flutter pub get
flutter run --dart-define=APP_MODE=passenger
flutter run --dart-define=APP_MODE=driver
flutter run -d chrome --dart-define=APP_MODE=admin
\end{code}
无可用 Chrome 时，已配置 Web 开发环境可选择 \texttt{flutter run -d web-server --dart-define=APP\_MODE=admin}，随后用浏览器访问命令给出的地址；不能把开发服务器当作真实业务后端。

\begin{longtable}{p{3cm}p{3cm}p{2.5cm}p{6.2cm}}
\toprule 入口 & 手机号 & 密码 & 用户 \\ \midrule
passenger & 13800000001 & 123456 & 张同学，初始信用 96 \\
driver & 13800000002 & 123456 & 王师傅 \\
admin & 13800000003 & 123456 & 李管理员 \\
\bottomrule
\end{longtable}
这些账号只用于课程原型，不是真实凭据。页面快捷按钮先填充账号，再提交登录；账户角色需匹配当前入口。

\section{高价值源码索引}
\begin{longtable}{p{7.6cm}p{7.1cm}}
\toprule 路径 & 阅读目的 \\ \midrule
\pathcode{lib/main.dart}、\pathcode{lib/app.dart} & 入口模式注入及 MaterialApp.router \\
\pathcode{lib/core/router/app_router.dart} & 登录鉴权和角色重定向 \\
\pathcode{lib/core/providers/auth_provider.dart} & 认证状态、退出、信用变更 \\
\pathcode{lib/core/persistence/} & 演示快照、本地存储、同源标签页同步、内存回退 \\
\pathcode{lib/features/auth/data/mock_auth_repository.dart} & 演示账号、密码及角色校验 \\
\pathcode{lib/features/passenger/booking/providers/passenger_booking_provider.dart} & 余座、创建预约、取消、支付 \\
\pathcode{lib/features/driver/tasks/providers/driver_task_provider.dart} & activeTask、统计周期、位置 Timer \\
\pathcode{lib/features/admin/providers/admin_provider.dart} & CRUD、绑定、调度确认、后台模拟位置 \\
\pathcode{lib/core/providers/notification_provider.dart} & 本地通知状态 \\
\pathcode{lib/features/*/presentation/} & 各端工作台与交互按钮 \\
\pathcode{lib/shared/widgets/} & 入口切换、通用 Scaffold 与卡片 \\
\bottomrule
\end{longtable}
核心模型为手写 Dart 类和 copyWith，没有代码生成前置步骤。依赖安装不需要运行 build\_runner，也不需要先启动 json-server 或 Supabase。

\section{演示顺序}
\begin{enumerate}
\item 乘客：快捷账号登录，预约大学城早班线 19 座，查看订单并模拟支付；另一次演示取消预约，核对信用 96→91 及通知。
\item 司机：查看乘客名单，切换日周月统计，开始 task-001，手动模拟一次位置记录，完成任务，检查按钮禁用及共享停止。
\item 后台：编辑 bus-002，新增车辆与司机，生成智能建议并确认，点击 bus-003 标记，查看固定客户 ETA，刷新模拟位置，再切换周统计。
\end{enumerate}
网页演示的数据保存在当前浏览器，同源标签页可看到对应业务域更新；刷新后重新登录即可恢复业务样例。不同浏览器和设备不共享数据，后台调度也不会创建其他角色的业务对象。需要重复演示时，先退出到登录页，确认“重置演示数据”；原生平台仍只保存本次会话状态。

\chapter{测试计划、用例与报告}
\toolmark{Flutter Test、WidgetTester、ProviderContainer、测试执行记录}
\section{测试目标、范围和策略}
测试覆盖三端主流程、业务拒绝路径、司机状态 getter / copyWith、Controller、模型与枚举。Widget 测试用 ProviderScope 注入指定入口，按 Key 定位按钮；业务单测用独立 ProviderContainer，完成后 dispose。当前自动化测试使用 Mock，不需要数据库、支付账号、GPS 设备或外部服务。

登录模拟 450ms，支付模拟 500ms；测试用 \texttt{tester.pump(Duration(...))} 推进假时钟。周期 Timer 及持续动画要显式控制，不能将 \texttt{pumpAndSettle} 当作等待所有后台任务的通用方法。滚动区域中的按钮先 \texttt{ensureVisible} 或 \texttt{scrollUntilVisible}；当前快捷登录 Key 是 \texttt{quick\_\$\{mode.value\}\_login}，不是旧测试样例的 mock 前缀。

\section{乘客用例}
\begin{longtable}{p{1.7cm}p{4.2cm}p{8.8cm}}
\toprule 编号 & 场景 & 操作及预期 \\ \midrule
TC-P-01 & 登录工作台 & passenger 入口，账号 13800000001；显示预约车次和我的预约 \\
TC-P-02 & 选座预约 & trip-001 选择 19 座，确认；生成订单并显示成功提示 \\
TC-P-03 & 取消与扣分 & 对新建待发车预约确认取消；扣 5 分，信用 96→91 \\
TC-P-04 & 模拟支付 & 创建新订单，选择微信支付并确认；paid 和本地成功通知 \\
TC-P-05 & 信用和通知 & 查看面板，预约 / 支付 / 取消后检查消息；部分由上述 Widget 合并覆盖 \\
TC-P-06 & 错误角色入口 & 乘客账号登录 driver；抛出当前入口不匹配异常 \\
TC-P-07 & 重复预约 & 同用户 trip-001 先预约 19 再预约 20；第二次拒绝 \\
TC-P-08 & 非 pending 取消 & booking-legacy-001/002；两者均提示只有待发车可取消 \\
\bottomrule
\end{longtable}
上述是需求场景，不等于八个独立自动化测试。乘客自动化测试分别位于 \pathcode{widget_test.dart} 和 \pathcode{business_rule_test.dart}；筛选切换、满座、弱网、真实支付和移动端兼容性仍需补充专项验证。

\section{司机专项用例和测试数据}
\begin{longtable}{p{4.6cm}r p{8.5cm}}
\toprule 测试组 & 实际项数 & 验证点 \\ \midrule
登录与工作台 & 7 & 标题、司机姓名、任务和乘客数量、初始位置状态、路线、暂无上报 \\
统计看板 & 4 & 默认日数据、周切换、月切换、准点率进度条 \\
任务列表 & 8 & 路线与车次号、起终点、车牌、人次、里程、标签、名单、上车点 \\
开始任务与上报 & 6 & 开始任务、首条位置、手动上报、坐标变化、车辆 ID、速度 \\
完成任务 & 3 & 完成状态、停止上报、提示消息 \\
边界情况 & 5 & 无任务时按钮禁用、单 active 任务、已完成不可操作、可开始下一任务 \\
State 单元测试 & 9 & 任务数、乘客数、activeTask、统计选择、最新位置、copyWith 和清空状态 \\
Controller 单元测试 & 7 & 启动、重复启动、不存在任务、完成、空任务完成、周期、上报数量 \\
模型与枚举 & 8 & 任务 copyWith、三种状态、乘客、位置、统计、结果字段 \\
合计 & 57 & Widget 33 项，单元测试 24 项 \\
\bottomrule
\end{longtable}
原司机报告虽然总数 57 正确，但“开始任务与位置上报 7 个”“Widget 34 个”与明细不符，本表已按实际测试源码更正。

\begin{longtable}{p{2cm}p{3.1cm}p{5.5cm}p{1.2cm}p{2.6cm}}
\toprule 任务 & 路线 & 起止站点 & 乘客 & 里程 \\ \midrule
task-001 & 大学城早班线 & 南门学生公寓→综合教学楼 & 3 & 8.6km \\
task-002 & 科技园通勤线 & 图书馆广场→科技园东门 & 4 & 12.4km \\
task-003 & 晚间返校线 & 附属医院站→北区宿舍 & 2 & 15.8km \\
\bottomrule
\end{longtable}
司机为王师傅，车辆 bus-001，车牌校A·1028；车次号依次为 D-20260501-01、02、03。模拟统计：日 4 趟 / 38 人 / 42.6km / 98\%；周 24 趟 / 236 人 / 286.4km / 95\%；月 98 趟 / 1034 人 / 1198.5km / 92\%。

\section{后台与负向业务用例}
\begin{longtable}{p{2cm}p{4.2cm}p{8.5cm}}
\toprule 编号 & 场景 & 预期或验证点 \\ \midrule
TC-A-01 & 管理员登录 & 后台工作台、车辆与司机面板 \\
TC-A-02 & 车辆编辑和新增 & 编辑 bus-002 车牌为校A·2040；新增校A·5098 显示在列表 \\
TC-A-03 & 司机新增 & 姓名、手机、驾照保存后显示孙师傅 \\
TC-A-04 & 建议生成与确认 & 显示 AI 推荐，确认后显示已确认 \\
TC-A-05 & 车辆监控 & 点 bus-003；显示校A·3099、钱师傅和位置详情 \\
TC-A-06 & ETA 与位置刷新 & 固定预计 8 分钟样例；手动模拟刷新仍正常展示 \\
TC-A-07 & 运营统计 & 切换周，显示样例人次 1298 及路线排行 \\
TC-N-01 & 删除运行车辆 & bus-001 删除失败：运行中车辆不可删除 \\
TC-N-02 & 车辆绑定冲突 & 新司机尝试绑定 bus-002，失败：车辆已绑定其他司机 \\
TC-N-03 & 人工调度容量不足 & 注入 10 座车辆及 36 人需求，失败：车辆座位数不足 \\
TC-N-04 & 无可用智能车辆 & 注入空车辆 / 司机列表，返回没有可用于调度的空闲车辆和司机 \\
\bottomrule
\end{longtable}
后台主流程合并于一个 Widget 测试，并非七项独立自动化测试。七个负向自动化测试包含角色错误、重复预约、非法取消以及上表四项后台拒绝规则；测试仓库覆盖不足和无车数据通过 override 注入。

\section{可复现的测试代码和命令}
\begin{code}
final container = ProviderContainer();
addTearDown(container.dispose);
final controller = container.read(driverTaskProvider.notifier);
expect(controller.startTask('task-001').success, isTrue);
expect(container.read(driverTaskProvider).activeTaskId, 'task-001');
expect(controller.completeActiveTask().success, isTrue);
expect(container.read(driverTaskProvider).activeTaskId, isNull);
\end{code}
\begin{code}
flutter test
flutter test test/driver_test.dart
flutter test test/business_rule_test.dart
flutter test test/widget_test.dart --plain-name \
  'passenger can pay booking and receive notification'
# Optional coverage collection; no existing 80% claim is assumed:
flutter test --coverage
\end{code}

\section{历史执行记录：2026-10-10}
本节保留文档整合时的历史验证记录，不代表 2026 年 11 月修订后的测试数量。2026-10-10 在 Linux、Flutter 3.47.5 stable、Dart 3.13.4 下执行源码验证。该次 SDK 安装依赖时自动更新了六项测试/传递依赖和分析配置；原记录说明这些改写未纳入当时文档交付。
\begin{longtable}{p{7.7cm}p{7cm}}
\toprule 执行内容 & 2026-10-10 历史结果 \\ \midrule
\texttt{dart format --output=none --set-exit-if-changed lib test} & 退出码 0，无需改写源码格式 \\
\texttt{flutter analyze} & 退出码 0 \\
\texttt{flutter test --reporter expanded} & 退出码 0，69 项全部通过 \\
\texttt{flutter build web --release --base-href /BUS-FLUTTER/ --dart-define=APP\_MODE=admin} & 退出码 0，生成后台 Web 构建 \\
\bottomrule
\end{longtable}
当时 69 项组成：\pathcode{driver_test.dart} 57 项、\pathcode{widget_test.dart} 5 项、\pathcode{business_rule_test.dart} 7 项；该次仅整合文档，没有修改业务源码。未执行真实网络后端、真实付款、真机 GPS、生产压测、安全测试、Windows 或移动平台发布测试。

\section{2026-11-08 修订与回归范围}
本轮本地验证使用 Flutter 3.47.7 / Dart 3.13.5；CI 固定 Flutter 3.47.9，与修订前最近一次线上构建版本一致。新增 \pathcode{auth_regression_test.dart}、\pathcode{payment_regression_test.dart}、\pathcode{admin_dispatch_regression_test.dart} 等针对缺陷与持久化的回归用例，检查跨入口快捷账号、稳定路由、取消后重登信用分、迟到的登录回调、重复模拟付款、付款中取消、退款账本、司机绑定冲突、资源变更后确认、浏览器快照恢复与重置等行为。司机退出后的 Timer 停止与缓存上限也属于本轮验证范围。

浏览器验证通过真实页面控件操作执行，区分“代码已增加规则”与“对应页面路径已经实际通过”。模拟支付和本地同步的成功不表示真实网关、数据库或跨设备服务已经接入。历史截图继续作为原有课程证据，不替代本轮回归结果。
本轮执行非改写格式检查、静态分析、完整自动化回归和后台入口 Web release 构建，均成功；自动化测试由历史 69 项增至 131 项。新增用例覆盖上述状态迁移、持久化恢复、跨标签页旧状态写入、旧编辑表单冲突，以及退出后的定位计时器行为。Overleaf 归档同步检查通过，共 40 项输入。本地环境没有安装 XeLaTeX/latexmk，因此本轮未重新编译 PDF，原 PDF 继续作为历史导出。

发布前的本地 Chromium 页面回归确认：单次快捷账号可跨入口登录，预约、模拟付款、取消与信用分会同步到另一个已打开标签页，重登和刷新保留业务数据，重置会登出其他标签页并恢复初始分数；司机退出期间停止位置上报，重登恢复，完成任务后停止且完成状态在刷新后保留。此处记录浏览器演示能力，不表示跨设备后台已经实现。

\section{历史测试结论与证据边界}
原司机报告记录 2026-06-12、Windows 11、57 项通过、约 5 秒。原乘客报告日期 2026-06-10，并引用 2026-05-28 的“12 个场景通过”旧记录；它不是目前完整测试数。原后台手工报告环境为 Windows 11、Edge、Flutter 3.44.1 / Dart 3.12.1、Web Server、Mock，记载核心手工流程通过且无阻断缺陷；本次不把该结论作为新一轮浏览器手测。

原后台报告另含 JUnit 5 / Maven / Java UserTest 示例及截图，仓库不存在 pom.xml、对应 Java 模型或该测试源。这五个测试的创建用户、信用更新、余额更新、创建预约、取消预约仅作为历史补充演示，不计入任何当前 Flutter 自动化统计，也不说明 Flutter 已具有余额业务。
\reportfigure{report-assets/admin/admin_image1.png}{历史 Java/JUnit 补充演示截图：不属于当前 Flutter 自动化测试}
以下为原后台报告中的手工截图，时间和数据来自原材料，仅用于保留界面流程证据：
\reportfigure{report-assets/admin/admin_image2.png}{历史后台手测：新增车辆对话框}
\reportfigure{report-assets/admin/admin_image3.png}{历史后台手测：新增车辆结果}
\reportfigure{report-assets/admin/admin_image4.png}{历史后台手测：车辆编辑对话框}
\reportfigure{report-assets/admin/admin_image5.png}{历史后台手测：车辆编辑结果}
\reportfigure{report-assets/admin/admin_image6.png}{历史后台手测：司机新增与绑定结果}
\reportfigure{report-assets/admin/admin_image7.png}{历史后台手测：智能建议生成、人工确认入口}
\reportfigure{report-assets/admin/admin_image8.png}{历史后台手测：调度条件不足的提示}
\reportfigure{report-assets/admin/admin_image9.png}{历史后台手测：自绘位置面板与固定客户 ETA}
\reportfigure{report-assets/admin/admin_image10.png}{历史后台手测：位置变化局部截图}
\reportfigure{report-assets/admin/admin_image11.png}{历史后台手测：运营分析界面之一}
\reportfigure{report-assets/admin/admin_image12.png}{历史后台手测：运营分析界面之二}

\section{剩余测试风险}
本轮新增了重复付款、付款中取消、调度重新校验和存储恢复等回归，但仍需扩大边界座位与满座、统计筛选、错误账户归属、存储配额与损坏、多个标签页同时修改和长时间运行的验证。位置与通知缓存已限制为 200 条，不等于完成内存性能测试。生产版本还需跨客户端集成、网络异常恢复、数据库权限、真实支付回调、性能、国际化、无障碍及移动端兼容性测试。通过率不等于覆盖率，更不等于生产可靠性结论。

\chapter{配置管理、CI/CD 与文档交付}
\toolmark{Git、GitHub、GitHub Actions、pubspec.lock、SQL 迁移}
\section{配置项与版本基线}
\begin{longtable}{p{4cm}p{10.7cm}}
\toprule 配置项 & 当前路径与管理方式 \\ \midrule
源码和自动化测试 & lib/、test/，Git 跟踪 \\
依赖和静态分析 & pubspec.yaml、pubspec.lock、analysis\_options.yaml \\
平台工程 & android/、ios/、web/、linux/、macos/、windows/ \\
CI/CD & .github/workflows/flutter-ci-cd.yml \\
预备数据库 & supabase/migrations/、supabase/seed/ \\
唯一文档正文 & main.tex；历史图像及原始证据在 report-assets/ \\
Overleaf 交付 & overleaf.zip，由 tools/build\_overleaf.py 重建 \\
入口与代理约定 & README.md 只保留链接；AGENTS.md 仅记录代理工作约定 \\
\bottomrule
\end{longtable}
build/、.dart\_tool/ 等本地产物不纳入版本控制。当前 \texttt{.idea/} 被忽略，\texttt{.vscode/} 忽略规则是注释，不应沿用原配置报告“所有 IDE 配置都已忽略”的笼统说法。真实支付密钥、服务端密钥和凭据不应写入公开文档；当前 App 并未读取原报告中列出的 Supabase 环境参数。

历史团队约定以 main 为课程主线，feature/*、fix/*、docs/* 为建议分支；提交类型 feat、fix、docs、test、refactor、chore。原文 v0.1--v0.5 对应从骨架到综合功能的阶段命名，是历史版本方案，不保证仓库已有对应 tag。当前 manifest 版本为 1.0.0+1。

核心预约、支付、调度、位置流程变更应先更新本文件中的需求、设计与风险，再实现并补测试；测试通过后经评审合并，并记录新基线。需求、设计、测试和交付基线如今分别对应本文件章节，而不引用已移除的分散报告路径。

\section{CI/CD 的实际执行顺序}
工作流在 main 的 push、指向 main 的 PR 及手动触发时运行，Ubuntu 环境安装固定的 Flutter 3.47.9 stable，然后按下列顺序验证：
\begin{code}
flutter pub get
dart format --output=none --set-exit-if-changed lib test
python3 tools/build_overleaf.py --check
flutter analyze
flutter test
flutter build web --release --base-href /BUS-FLUTTER/ \
  --dart-define=APP_MODE=admin
\end{code}
CI 与本地使用相同的非改写格式门禁，格式不符即失败；Overleaf 检查同时阻止正文或素材修改后忘记更新 ZIP。工作流并不传 \texttt{BACKEND=mock}，因为当前不存在 BACKEND 读取逻辑。

只有 push 到 main 时上传 \pathcode{build/web} Pages artifact 并执行部署；PR 只验证。GitHub Pages 需选择 GitHub Actions 作为构建与部署来源；Web 子路径固定使用 \texttt{/BUS-FLUTTER/}。去掉 base-href 会使子路径静态资源加载错误。

\section{预备 Supabase SQL 与初始化顺序}
\pathcode{supabase/migrations/202605110001_initial_schema.sql} 包含 profiles、vehicles、drivers、trips、bookings、payments、dispatch\_demands、dispatch\_plans、vehicle\_locations、notifications，定义状态枚举、约束、索引、更新时间触发器、辅助权限函数和 RLS policy。
活动预约的座位和用户车次组合有条件唯一索引；driver 的车辆绑定列 unique；profiles 与 Supabase \texttt{auth.users} 外键相连。数据库 payment\_method 还有 manual，而当前 Dart 只含 wechat/alipay，未来接入需显式适配。

若专门验证数据库，在隔离开发 Supabase 实例中先执行迁移，随后创建三名 Auth 用户，最后执行种子文件：
\begin{code}
13800000001@school-bus.local / 123456
13800000002@school-bus.local / 123456
13800000003@school-bus.local / 123456
supabase/seed/202605110001_seed_demo_data.sql
\end{code}
种子依赖这三名 Auth 用户的 ID，未先创建用户会缺少相应 profile。当前仓库没有 Supabase CLI 配置、客户端初始化或后端切换，不编造 \texttt{supabase start} / 自动 reset 的现成工作流。现有 RLS 允许用户修改某些自身 profile/payment 字段，不能据此认为服务端信用和付款状态已被安全保护；生产接入需专项权限审查与事务设计。

\section{单一文档打包与维护}
\texttt{main.tex} 为唯一维护的正文，\texttt{report-assets/uml/} 为 UML 与过程图片，\texttt{plan/} 为原计划图，\texttt{admin/} 为后台历史截图，\texttt{evidence/} 为 StarUML、MPP、CSV 可编辑证据。2026 年 11 月从仓库已有的 \pathcode{docs/overleaf.zip} 恢复这套综合源文件，并保留用户上传的 docs 历史资料；后续只更新根目录正文和根目录交付包。新 ZIP 顶层直接包含 main.tex 与 .latexmkrc，没有外层项目目录，没有旧 Word/PDF/Markdown 报告，也不要求 Overleaf 访问仓库外部路径。

\begin{code}
python3 tools/build_overleaf.py
python3 tools/build_overleaf.py --check
# When a local TeX Live / XeLaTeX toolchain is installed:
latexmk -xelatex -outdir=build/report main.tex
\end{code}
脚本仅使用 Python 标准库，检查所有图像已被正文引用、资源存在、正文没有拆分输入，并以固定 ZIP 元数据打包。check 模式验证 ZIP 文件列表、CRC 和每项内容与当前源文件一致。修改正文、图片或证据后需重建交付包，不能只改 ZIP 或恢复多份平行正文。

\chapter{后续全栈技术方案}
\section{保留的原始选型及其状态}
本章整合最初技术方案，作为从课程原型走向真实系统的路线图。初稿 Riverpod 2.x、GoRouter 14 等依赖版本已经不代表当前项目；正式新增依赖时应按实际 Flutter / Dart 版本重新评估，而不是直接复制旧 pubspec 清单。
\begin{longtable}{p{3cm}p{5.2cm}p{6.5cm}}
\toprule 层 & 原方案候选 & 后续接入工作 \\ \midrule
网络 & Dio + Retrofit & 定义仓库接口、错误映射、超时重试、鉴权、真实 API \\
本地存储 & flutter\_secure\_storage、Hive & Token 安全存储、缓存、失效和账号隔离 \\
实时通信 & WebSocket / socket\_io\_client & 鉴权、重连、订阅、消息校验和序号 \\
定位与地图 & geolocator、flutter\_map + OSM / 高德 & 平台权限、坐标系、后台采集和可视化 \\
支付 & 微信 / 支付宝 SDK，初稿提及 pay & 核对官方能力；pay 不等于两种支付的现成后端方案 \\
推送 & firebase\_messaging / FCM & 平台配置、推送 token、服务端通知与隐私 \\
后端 & json-server 原型 / Supabase & 持久化、认证、RLS、服务端业务事务 \\
图表 & fl\_chart & 替换演示图表并连接真实统计数据 \\
数据生成 & json\_serializable / freezed & 只有选择生成模型后才引入生成命令 \\
调度辅助 & 本地规则 / OpenAI & 保留人工确认、安全校验与可解释约束 \\
\bottomrule
\end{longtable}
选择 Supabase 时可使用预备 SQL，但必须实现 Flutter 客户端和真实数据仓库，不能通过一个未被读取的 dart-define 开关声称接入完成。json-server 仅可作为另行搭建的接口样机，仓库当前也没有 db.json。

\section{目标 REST API}
以下是接口设计，不是可请求的现有服务。当前操作直接调用 Notifier。生产实现需定义请求字段、响应结构、错误码、角色权限、分页、幂等键及并发事务。
\begin{code}
POST   /api/auth/login
POST   /api/auth/register
GET    /api/trips
POST   /api/bookings
DELETE /api/bookings/:id
GET    /api/bookings/mine
POST   /api/payments
GET    /api/driver/tasks
GET    /api/driver/stats
GET    /api/admin/vehicles
POST   /api/admin/vehicles
PUT    /api/admin/vehicles/:id
DELETE /api/admin/vehicles/:id
GET    /api/admin/dispatch/suggest
POST   /api/admin/dispatch/confirm
GET    /api/admin/analytics
WS     /ws/location
\end{code}
司机 CRUD、订单支付查询、回调、退款、通知与统计导出尚需补齐接口契约。模拟订单、任务、后台需求 ID 与数据库 UUID 的映射也要明确，不能直接把当前演示字符串当数据库主键。

\section{实时位置、推送与 ETA}
原方案设计司机每 5 秒发出 \texttt{driver:location}，包含 vehicleId、lat、lng、timestamp，后台订阅 \texttt{vehicles:locationUpdate}。Socket.IO 事件协议与裸 WebSocket 路由不是同一种接口约定，服务端与客户端需统一后再接入。
\begin{code}
// Planned Socket.IO event shape; not existing application code.
socket.emit('driver:location', {
  'vehicleId': vehicleId,
  'lat': latitude,
  'lng': longitude,
  'timestamp': timestamp,
});
socket.on('vehicles:locationUpdate', handleLocationUpdate);
\end{code}
正式客户端发车后启动 GPS、完成后停止采集，并在断线时恢复正确订阅；位置数据应有授权、精度、时效、保留期限与异常处理。乘客定位必须取得明确同意。ETA 可先用 Haversine 直线距离和假设速度做估算，再替换为路网与交通预测；必须标注估算误差，不得把当前固定 8/14/23 分钟样例当定位计算结果。

\section{支付、数据一致性与调度演进}
真实预约需在服务器以事务保护余座与座位唯一性；真实支付必须由服务端签名下单、验签回调与对账，客户端只能发起和展示状态，不能自行把订单改为已付。取消应统一计算信用、退款与座位释放，涉及支付中的竞态时使用幂等机制。
智能调度可从当前贪心规则演进到路线聚合、多车分拆、司机排班、车辆维护窗口及全局优化；每条推荐应说明容量、负载和约束，人工确认时再次校验资源可用性。如果使用 OpenAI，仅将其作为辅助，不应绕过确定性安全约束或上传未经批准的个人位置数据。

\chapter{工具证据、整合追踪与总结}
\section{课程工具证据索引}
\begin{longtable}{p{2.5cm}p{4.5cm}p{7.7cm}}
\toprule 阶段 & 工具标识 & 当前证据 \\ \midrule
计划 & Project 类工具、CSV、甘特图 & 第二章，plan 图像，evidence/project\_plan.mpp 与 project\_plan.csv \\
需求 & StarUML、规格说明与追踪矩阵 & 第三章，用例 PNG，StarUML 原模型 \\
设计 & StarUML、架构与平台设计 & 第四至六章，12 张 UML 和 8 张过程截图 \\
源码 & Flutter/Dart、Riverpod、GoRouter & 第七章，lib/、pubspec.yaml \\
测试 & Flutter Test、WidgetTester、ProviderContainer & 第八章，test/、历史后台截图、本次实测结果 \\
配置 & Git、锁文件、GitHub Actions、迁移 SQL & 第九章，工作流、pubspec.lock、supabase/ \\
整合交付 & LaTeX、XeLaTeX、Python ZIP & main.tex、.latexmkrc、tools/build\_overleaf.py、overleaf.zip \\
\bottomrule
\end{longtable}
编辑器和原生建模工具的使用来自历史记录与截图，本文不声称本次重新在 Project 或 StarUML 中操作过所有模型。

\section{原材料合并与更正记录}
\begin{longtable}{p{4cm}p{5cm}p{5.7cm}}
\toprule 原材料类别 & 整合位置 & 去重 / 更正说明 \\ \midrule
根目录 README 与技术方案 & 概述、运行、扩展方案 & 旧框架版本及全栈服务为规划，不作为当前依赖 \\
计划 Word 与 PDF & 第二章 & 相同正文合并一次，提取三张图片；6 月 4 日与 25 日分别说明 \\
需求 Markdown & 第三至五章 & 合并模型字典、用例、状态、追踪；修正 UML 数量与枚举 \\
Windows 设计 DOC 与 PDF & 第六章 & 保留原生宿主设计，区分模板能力与未实现业务 \\
源码说明 Markdown & 第四、七章 & 去掉原作者本机绝对路径，以仓库相对路径解释 \\
司机测试 Word 与 PDF & 第八章 & 三部分合并；实际 57 项，修正分组数量与旧 Key \\
乘客测试 Word 与 PDF & 第八章 & 保留场景与风险，去掉已接入 Supabase 的错误说法 \\
后台测试 Word 与 PDF & 第八章 & 提取 12 张截图，JUnit 仅作为历史补充，按 Flutter 源码说明实测 \\
配置管理 PDF & 第九章 & 保留版本基线；纠正 BACKEND、IDE 忽略规则和 CI 命令 \\
工具标识、最终报告 & 本章及各阶段 & 不再引用缺失的拆分报告、PlantUML 文件或旧交付目录 \\
StarUML 过程说明及图片 & 第五章 & 保留原模型、全部 20 张 PNG；说明早期设计与当前实现差异 \\
平台启动图 README & 第六章 & 说明融入平台章节，不再保留额外正文 \\
MPP、CSV、MDJ & report-assets/evidence/ & 原始可编辑证据保留，不产生平行文档 \\
\bottomrule
\end{longtable}
Word 与其 PDF 导出版本内容相同的部分不重复收录；历史截图未经伪造或改写。原报告引用但仓库实际不存在的独立司机测试计划/样例、配置 Word、PlantUML 等文件，不再作为已交付证据列出。

\section{最终交付与结论}
唯一阅读正文为本综合报告；Overleaf 项目由 main.tex、编译配置、图像和可编辑计划/建模证据组成，导入后编译得到一份完整 PDF。源码、测试、平台工程和预备 SQL 仍作为工程资产保留，不被文档合并替代。

当前课程原型具备三类工作台与网页本地保存、同源标签页同步能力，本轮强化了快捷登录、信用分延续、模拟付款状态一致性、调度确认及司机定位生命周期。历史验证结果和本轮回归范围在测试章节分别记载。正式上线仍需要真实认证、服务端持久化、跨设备业务一致性、实际支付与定位、权限审查以及生产级测试；当前浏览器本地保存不承担这些职责。
\end{document}
